\section{Thinking in Shapes: Flat, Round, and Saddle}

To understand the crystal, you first need a feel for the kind of space it lives in. Space comes in three basic flavors, and you already know all three from everyday surfaces. The trick is to stop thinking of them as surfaces sitting in our world and start thinking of them as whole worlds in their own right.

\subsection{The three flavors}

Take a flat tabletop. Draw a triangle. Its three angles add up to exactly 180 degrees. Parallel lines stay the same distance apart forever. A circle of radius two has exactly twice the edge of a circle of radius one. This is \textbf{flat} space, the geometry of Euclid, the one we learn in school.

Now take the surface of a ball, like the Earth. Draw a big triangle, say from the North Pole down to the equator, along the equator a quarter of the way around, and back up. Every corner is a right angle. The angles add up to 270 degrees, more than 180. Lines that start out parallel, like two lines of longitude at the equator, bend toward each other and meet at the pole. Room runs out. This is \textbf{round}, or positively curved, space. A sphere is finite. You can walk all the way around it and come home.

Now imagine the opposite. A surface shaped like a saddle, or the ruffled edge of a lettuce leaf, or a Pringle. Draw a triangle on it and the angles add up to less than 180 degrees. Lines that start parallel curve away from each other and never meet. There is always more room than you expect. This is \textbf{saddle}, or negatively curved, space. Its proper name is \textbf{hyperbolic} space, and it is the star of this walkthrough.

\subsection{How to tell which world you are in, from inside}

Here is the key move. You do not need to step outside a space to know its curvature. You can measure it from within, by checking how much room there is.

Stand at a point and ask: how many places are within one step of me? Within two steps? Within ten?

\begin{itemize}
\item In flat space, the amount of room grows steadily. The area inside a circle grows with the square of the radius. Double the radius, four times the room.
\item In round space, the room grows and then shrinks back as you wrap around. It runs out.
\item In hyperbolic space, the room grows faster and faster, exploding as you go out. The amount within a given distance does not just grow, it grows exponentially.
\end{itemize}

That last fact is the whole reason hyperbolic space matters for us. It has, in a precise sense, endless room near every point. The further you look, the more the world opens up.

\subsection{Why this matters for the crystal}

The crystal of experience has to grow forever, adding more and more cells. Flat space does not have enough room for that kind of explosive, branching growth. It fills up too slowly. Round space runs out entirely. Only hyperbolic space, with its room that keeps opening up, can comfortably hold a structure that branches and grows without end.

So the very first thing the theory needs from geometry is negative curvature. Not because someone liked saddles, but because the universe grows, and only a saddle-shaped world has the room.

\subsection{Why our space does not look curved}

If space is really saddle-shaped, why does the room around you look flat, with parallel lines staying parallel? For the same reason the Earth looks flat from a sidewalk. Any curved space looks flat in a small enough patch, and the curvature only shows over large distances or in fine structure. We live in a patch, so day to day space looks ordinary and Euclidean, exactly as it should. The hyperbolic character of the substrate reveals itself not in your living room but in the deep things it buys, the endless room to grow and the absence of any preferred direction. Curved underneath, flat to the touch.

\textbf{Takeaway:} space can be flat, round, or saddle-shaped, and you can tell which from inside by how fast room grows. Hyperbolic (saddle) space has room that opens up exponentially, which is exactly what a forever-growing crystal needs.
